% ASSERT_CONVENTION: entropy_base=nats, generator_convention=probabilist_dpdt_pQ, matrix_norm=sup_norm_rows, experiential_density=I*(1-I/H)
%
% Theorem A Lemma Assembly -- Plan 01-01
% Phase: 01-theorem-a-assembly
%
% All expressions use nats (ln), probabilist generator convention
% (row-vector left: dp/dt = pQ, piQ = 0), and sup-norm on rows.
%
% References:
%   [BEGK02]  Bovier, Eckhoff, Gayrard, Klein, Comm. Math. Phys. 228 (2002) 219--255
%   [BdH15]   Bovier, den Hollander, Metastability: A Potential-Theoretic Approach, Springer (2015)
%   [FW12]    Freidlin, Wentzell, Random Perturbations of Dynamical Systems, 3rd ed., Springer (2012)
%   [CV16]    Champagnat, Villemonais, Probab. Theory Relat. Fields 164 (2016) 243--283
%   [DV75-83] Donsker, Varadhan, Comm. Pure Appl. Math. (1975--1983)

\documentclass[11pt]{article}
\usepackage{amsmath,amssymb,amsthm}
\usepackage[margin=1in]{geometry}

\newtheorem{lemma}{Lemma}
\newtheorem{theorem}{Theorem}
\newtheorem{remark}{Remark}
\newtheorem{definition}{Definition}

\newcommand{\eps}{\varepsilon}
\newcommand{\Ds}{\Delta_s}
\newcommand{\Db}{\Delta_b}
\newcommand{\Bstable}{B_{\mathrm{stable}}}
\newcommand{\BBB}{B_{\mathrm{BB}}}
\newcommand{\Teps}{T_\eps}
\newcommand{\TV}{\mathrm{TV}}
\newcommand{\QSD}{\mathrm{QSD}}

\title{Theorem A: Lemma Statements with Explicit Error Terms}
\author{Experiential Measure Formalization}
\date{}

\begin{document}
\maketitle

%======================================================================
\section{Setup and Notation}
%======================================================================

Let $(\Omega, Q_\eps)_{\eps > 0}$ be a family of finite irreducible
continuous-time Markov chains with state space
$\Omega = B \times M$ ($|B|, |M| < \infty$) and generators of
Metropolis type:
\[
  Q_\eps(x,y) = r(x,y) \exp\!\bigl(-[E(y) - E(x)]^+ / \eps\bigr),
  \qquad x \neq y,
\]
where $r(x,y) > 0$ for all pairs $(x,y)$ with $x \neq y$ (irreducibility),
$r(x,y) = r(y,x)$ (reversibility with respect to the Gibbs measure
$\pi_\eps(x) \propto \exp(-E(x)/\eps)$), and $[a]^+ = \max(a, 0)$.

Row sums are zero: $Q_\eps(x,x) = -\sum_{y \neq x} Q_\eps(x,y)$.
Convention: row-vector left multiplication, $dp/dt = pQ_\eps$.

The experiential density functional is
$\rho(p) = I(B;M) \cdot (1 - I(B;M)/H(B))$,
with entropy and mutual information computed in nats ($\ln$).
The trajectory functional is
$\mu([0,T]) = \int_0^T \rho(p_t)\, dt$ (units: nat-seconds).

\medskip
\textbf{Standing assumptions for Theorem~A:}
\begin{enumerate}
\item The Freidlin--Wentzell cycle hierarchy (Lemma~1) decomposes
  $\Omega$ into metastable sets including $\Bstable$ with
  communication height $\Ds$ and $\BBB$ with communication height
  $\Db$, satisfying $\Ds > \Db$.
\item \textbf{Density bounds:} $\rho(\nu_\QSD^{\Bstable}) \geq c > 0$
  (the quasi-stationary distribution on $\Bstable$ has non-trivial
  self-modeling), and $\rho(p) \leq \rho_{\max} = H(B)/4$ for all
  distributions $p$ on $\Omega$ (automatic for finite state spaces).
\item The observation horizon is
  $\Teps = \exp((\Ds - \alpha)/\eps)$ for a fixed
  $\alpha \in (0, \Ds - \Db)$.
\item The initial distribution $p_0$ is concentrated on $\Bstable$.
\end{enumerate}

%======================================================================
\section{Lemma Statements}
%======================================================================

%----------------------------------------------------------------------
\begin{lemma}[Basin Partition]\label{lem:partition}
\textbf{Source:} Freidlin--Wentzell \cite{FW12}, Chapter~6,
Theorem~6.3.1 (cycle decomposition of Markov chains).

\medskip
\textbf{Statement.}
For the family $(Q_\eps)_{\eps > 0}$ of reversible Metropolis generators
on the finite state space $\Omega$ with energy function $E$,
the Freidlin--Wentzell cycle hierarchy defines a partition of $\Omega$
into metastable sets (cycles) $\{B_1, B_2, \ldots, B_k\}$ at
each hierarchy level $h$, with communication heights
\[
  \Delta_i = V(B_i, B_i^c) - \min_{x \in B_i} E(x),
\]
where $V(A, A^c) = \min_{x \in A, y \notin A}
V(x,y)$ is the minimal communication height between $A$ and its
complement, defined via
\[
  V(x,y) = \min_{\gamma: x \to y} \max_{i=0,\ldots,|\gamma|-1}
  E(\gamma_i) - E(x)
\]
(minimum over all paths $\gamma$ from $x$ to $y$ of the maximal
energy barrier along the path, measured from the starting energy).

Identify $\Bstable$ and $\BBB$ as two cycles with
$\Ds > \Db > 0$.

\medskip
\textbf{Error term:} None.
This is an exact graph-theoretic construction on the finite state space.
No asymptotic approximation is involved.

\medskip
\textbf{Output type:}
Partition $(\Bstable, \BBB, \Omega \setminus (\Bstable \cup \BBB))$
with communication heights $\Ds, \Db \in \mathbb{R}_{>0}$.
\end{lemma}
%----------------------------------------------------------------------

\begin{remark}
The communication height $V(x,y)$ in Freidlin--Wentzell theory is
defined for the continuous-time chain via the W-graph (graph of
optimal transition paths). For finite state spaces with Metropolis
rates, it coincides with the energy barrier: $V(x,y) = \max_\gamma
E(\gamma_*) - E(x)$ along the optimal path. The partition is
determined entirely by the energy landscape $E$ and the adjacency
structure, independent of $\eps$.
\end{remark}

%----------------------------------------------------------------------
\begin{lemma}[Residence Time Lower Bound]\label{lem:residence}
\textbf{Source:} Bovier, Eckhoff, Gayrard, Klein \cite{BEGK02},
Theorem~1.2 (mean exit time via capacities); see also
Bovier--den~Hollander \cite{BdH15}, Chapter~7, Theorem~7.8.

\medskip
\textbf{Statement.}
Let $\tau_{\Bstable^c} = \inf\{t > 0 : X_t \notin \Bstable\}$ be
the first exit time from $\Bstable$. For any starting state
$x \in \Bstable$,
\begin{equation}\label{eq:mean-exit}
  \mathbb{E}_x[\tau_{\Bstable^c}]
  = \frac{1}{\pi_\eps(x) \cdot \mathrm{cap}_\eps(x, \Bstable^c)}
    \cdot (1 + \delta_2),
\end{equation}
where $\pi_\eps$ is the stationary (Gibbs) distribution,
$\mathrm{cap}_\eps(x, \Bstable^c)$ is the capacity between $x$
and $\Bstable^c$ (defined as
$\mathrm{cap}_\eps(x,A) = \sum_{y} \pi_\eps(x) Q_\eps(x,y)
(h_A(x) - h_A(y))$ with $h_A$ the equilibrium potential), and
\begin{equation}\label{eq:delta2}
  |\delta_2| \leq C_2 \exp(-c_2 / \eps)
\end{equation}
with $c_2 > 0$ determined by the sub-leading communication height
in the cycle hierarchy.

Specifically, from the capacity asymptotics in
\cite{BEGK02} Theorem~1.2 and \cite{BdH15} Theorem~7.8,
the capacity satisfies
$\mathrm{cap}_\eps(x, \Bstable^c) = \kappa_s \exp(-\Ds/\eps)
(1 + O(\exp(-c_2'/\eps)))$
where $\kappa_s > 0$ is a prefactor depending on the curvatures of
$E$ at the saddle point and the minimum, and $c_2' > 0$.
Since $\pi_\eps(x) = Z_\eps^{-1} \exp(-E(x)/\eps)$, combining gives
\begin{equation}\label{eq:mean-exit-sharp}
  \mathbb{E}_x[\tau_{\Bstable^c}]
  = K_s \cdot \exp(\Ds / \eps) \cdot (1 + \delta_2),
\end{equation}
where $K_s = Z_\eps / (\kappa_s \exp(-E(x)/\eps))$ is a prefactor
that is $O(1)$ in $\eps$ (polynomial in $|\Omega|$ and the rate
prefactors $r(x,y)$).

The rate constant $c_2$ satisfies
\begin{equation}\label{eq:c2-bound}
  c_2 \geq \min\bigl(\Ds - \Db,\; \Delta_{\mathrm{sub}}\bigr),
\end{equation}
where $\Delta_{\mathrm{sub}}$ is the sub-leading communication height
(the smallest gap between the leading exit barrier $\Ds$ and the next
barrier in the cycle hierarchy within $\Bstable$).

In particular, $c_2 > 0$ and $c_2 \geq \Ds - \Db > 0$ by the
standing assumption $\Ds > \Db$.

\medskip
\textbf{Exponential law} (\cite{BEGK02}, Theorem~1.4):
The rescaled exit time converges in distribution:
\begin{equation}\label{eq:exp-law}
  \frac{\tau_{\Bstable^c}}{\mathbb{E}_x[\tau_{\Bstable^c}]}
  \xrightarrow{d} \mathrm{Exp}(1)
  \qquad \text{as } \eps \to 0.
\end{equation}
This provides concentration of $\tau_{\Bstable^c}$ around its mean
without requiring full Donsker--Varadhan large deviation machinery.
Quantitatively, for any $\eta > 0$:
\begin{equation}\label{eq:exp-concentration}
  \mathbb{P}_x\!\left(
    \left|\frac{\tau_{\Bstable^c}}
              {\mathbb{E}_x[\tau_{\Bstable^c}]} - 1\right|
    > \eta
  \right)
  \leq 2\exp(-\eta/2)
  \qquad \text{for } \eps \text{ sufficiently small},
\end{equation}
using the exponential distribution tail bound
$\mathbb{P}(|X - 1| > \eta) \leq 2e^{-\eta/2}$ for $X \sim \mathrm{Exp}(1)$
(valid for $\eta < 1$; for $\eta \geq 1$, use $\mathbb{P}(X > 1+\eta)
= e^{-(1+\eta)}$).

\medskip
\textbf{Output type:}
Lower bound on $\mathbb{E}_x[\tau_{\Bstable^c}]$ of the form
$K_s \exp(\Ds/\eps)(1 + \delta_2)$ (an expectation with multiplicative
error), plus the exponential law providing distributional convergence
(hence concentration of the random variable $\tau_{\Bstable^c}$
around its mean).
\end{lemma}
%----------------------------------------------------------------------

%----------------------------------------------------------------------
\begin{lemma}[QSD Convergence]\label{lem:qsd}
\textbf{Source:} Champagnat--Villemonais \cite{CV16}, Theorem~2.1
(exponential convergence to QSD for absorbed Markov processes on
finite state spaces).

\medskip
\textbf{Statement.}
Consider the killed chain: the Markov chain $X_t$ restricted to
$\Bstable$ with absorption upon hitting $\Bstable^c$. Let
$Q_\eps^{\Bstable}$ denote the generator of this killed chain
(obtained by restricting $Q_\eps$ to $\Bstable$ and setting
absorption rates at the boundary).

There exists a unique quasi-stationary distribution
$\nu_\QSD$ on $\Bstable$ satisfying
\[
  \mathbb{P}_{\nu_\QSD}(X_t \in \cdot \mid t < \tau_{\Bstable^c})
  = \nu_\QSD(\cdot)
  \qquad \text{for all } t \geq 0.
\]
Moreover, for any initial distribution $p_0$ supported on $\Bstable$,
\begin{equation}\label{eq:qsd-convergence}
  \left\|\mathbb{P}_{p_0}(X_t \in \cdot \mid t < \tau_{\Bstable^c})
  - \nu_\QSD\right\|_{\TV}
  \leq C_3 \exp(-\gamma_D \cdot t),
\end{equation}
where $\gamma_D > 0$ is the spectral gap of the killed generator
$Q_\eps^{\Bstable}$ (specifically, the gap between the two
largest eigenvalues of $Q_\eps^{\Bstable}$, noting all eigenvalues
are negative), and $C_3 > 0$ depends on the initial distribution
$p_0$.

\medskip
\textbf{Key property of $\gamma_D$:}
The spectral gap $\gamma_D$ of the killed chain within $\Bstable$
is $O(1)$ as $\eps \to 0$ (not exponentially small). This holds
because:
\begin{itemize}
\item The killed chain operates within a single metastable basin
  $\Bstable$.
\item All transition rates \emph{within} $\Bstable$ are $O(1)$ or
  larger: internal transitions have barriers $\leq \Delta_{\mathrm{int}}
  < \Ds$, so rates $\geq r_{\min} \exp(-\Delta_{\mathrm{int}}/\eps)$.
  The exponentially small rates are the \emph{inter-basin} rates
  (barriers $\geq \Ds$), which are removed by the killing.
\item By the Cheeger inequality for the killed chain on $\Bstable$
  (see, e.g., \cite{BdH15} Section~7.2), $\gamma_D$ is bounded below
  by a function of the minimum within-basin transition rate and the
  basin geometry, both of which are $O(1)$ in $\eps$.
\end{itemize}

More precisely, $\gamma_D \geq \gamma_{\min}$ where $\gamma_{\min}$
depends on $|\Bstable|$, $r_{\min}$, and $\Delta_{\mathrm{int}}$
(the maximum internal barrier within $\Bstable$) but \emph{not}
on $\Ds$ or $\eps$ (beyond the within-basin rates).

\medskip
\textbf{Error term:}
\[
  |\delta_3(t)| = C_3 \exp(-\gamma_D t).
\]
This is not an asymptotic-in-$\eps$ error: it is an error
in time $t$ for fixed $\eps$. The relevant statement is that
after time $t_{\mathrm{mix}} = O(1/\gamma_D) = O(1)$ within
the basin, the conditional distribution is within $\TV$-distance
$C_3 e^{-\gamma_D t_{\mathrm{mix}}}$ of $\nu_\QSD$.

Since the residence time in $\Bstable$ is
$\tau \sim \exp(\Ds/\eps) \gg 1/\gamma_D$ for small $\eps$,
the chain spends all but an $O(1)$-duration transient at the QSD.

\medskip
\textbf{Output type:}
The unique QSD $\nu_\QSD$ on $\Bstable$, with a convergence rate
$\gamma_D = O(1)$ and the TV-distance bound \eqref{eq:qsd-convergence}
(a high-probability statement about the conditional distribution
within the basin).
\end{lemma}
%----------------------------------------------------------------------

%----------------------------------------------------------------------
\begin{lemma}[Stable Measure Lower Bound]\label{lem:stable-lower}
\textbf{Source:} Combination of Lemma~\ref{lem:residence}
(residence time) and Lemma~\ref{lem:qsd} (QSD convergence).

\medskip
\textbf{Statement.}
The experiential measure accumulated during residence in $\Bstable$
satisfies
\begin{equation}\label{eq:mu-stable}
  \mu_{\mathrm{stable}}
  = \int_0^{\tau_{\Bstable^c}} \rho(p_t)\, dt
  \geq c \cdot K_s \cdot \exp(\Ds/\eps) \cdot (1 - \delta_4),
\end{equation}
where $c = \rho(\nu_\QSD) > 0$ (the standing density-bound assumption),
$K_s$ is the residence-time prefactor from Lemma~\ref{lem:residence},
and
\begin{equation}\label{eq:delta4}
  \delta_4 = O\!\left(\frac{1}{\gamma_D \cdot K_s \exp(\Ds/\eps)}\right)
  = O(\exp(-\Ds/\eps)).
\end{equation}

\medskip
\textbf{Derivation of the bound.}
Decompose the residence period $[0, \tau_{\Bstable^c}]$ into
a transient phase $[0, t_0]$ and a QSD phase $[t_0, \tau_{\Bstable^c}]$
where $t_0 = O(1/\gamma_D)$ is chosen so that
$\|p_{t_0}^{\mathrm{killed}} - \nu_\QSD\|_\TV < \delta$ for a small
$\delta > 0$.

During $[t_0, \tau_{\Bstable^c}]$:
\[
  |\rho(p_t) - \rho(\nu_\QSD)|
  \leq L_\rho \|p_t - \nu_\QSD\|_\TV
  \leq L_\rho C_3 e^{-\gamma_D t}
\]
where $L_\rho$ is the Lipschitz constant of $\rho$ with respect to
TV distance (bounded by $O(H(B))$ on finite state spaces).
Therefore
\begin{align*}
  \mu_{\mathrm{stable}}
  &= \int_0^{t_0} \rho(p_t)\, dt
     + \int_{t_0}^{\tau_{\Bstable^c}} \rho(p_t)\, dt \\
  &\geq 0 + \int_{t_0}^{\tau_{\Bstable^c}}
    [\rho(\nu_\QSD) - L_\rho C_3 e^{-\gamma_D t}]\, dt \\
  &\geq c \cdot (\tau_{\Bstable^c} - t_0)
    - L_\rho C_3 / \gamma_D \\
  &= c \cdot \tau_{\Bstable^c}
    \Bigl(1 - \frac{t_0}{\tau_{\Bstable^c}}
    - \frac{L_\rho C_3}{c \gamma_D \tau_{\Bstable^c}}\Bigr).
\end{align*}
Since $t_0 = O(1)$ and
$\tau_{\Bstable^c} \sim K_s \exp(\Ds/\eps)$, both correction terms
are $O(\exp(-\Ds/\eps))$, giving $\delta_4 = O(\exp(-\Ds/\eps))$.

\medskip
\textbf{Rate constant:}
The error rate in $\delta_4$ is
\begin{equation}\label{eq:c4}
  c_4 = \Ds,
\end{equation}
which satisfies $c_4 = \Ds > \Ds - \Db > 0$.
The correction $\delta_4$ is superexponentially small compared to
the gap $\Ds - \Db$ that controls the final bound in Lemma~\ref{lem:ratio}.

\medskip
\textbf{Output type:}
Lower bound on $\mu_{\mathrm{stable}}$ with multiplicative error
$(1 - \delta_4)$ where $\delta_4 = O(\exp(-\Ds/\eps))$.
This is a statement about a random variable (the integrated
experiential measure during one residence period).

To convert to a high-probability bound: apply the exponential law
(Lemma~\ref{lem:residence}, eq.~\eqref{eq:exp-law}). Since
$\tau_{\Bstable^c} / \mathbb{E}[\tau_{\Bstable^c}] \to \mathrm{Exp}(1)$,
for any $\eta > 0$,
\[
  \mathbb{P}\bigl(\tau_{\Bstable^c}
  > (1-\eta) \mathbb{E}[\tau_{\Bstable^c}]\bigr)
  \geq 1 - (1 - e^{-(1-\eta)}) \to 1 - (1 - e^{-1}) \approx 0.632
\]
as $\eps \to 0$. Stronger: with probability $1 - e^{-a}$ the residence
time exceeds $a \cdot \mathbb{E}[\tau]$ (for $a < 1$). The
$\mu_{\mathrm{stable}}$ bound inherits this concentration.
\end{lemma}
%----------------------------------------------------------------------

%----------------------------------------------------------------------
\begin{lemma}[BB Occupation Upper Bound]\label{lem:bb-upper}
\textbf{Source:} Renewal theory for metastable chains; BEGK
\cite{BEGK02} Theorem~1.2 for exit times from $\BBB$; \cite{BdH15}
Chapter~7 for the excursion structure.

\medskip
\textbf{Statement.}
During the observation horizon $\Teps = \exp((\Ds - \alpha)/\eps)$,
the expected total occupation time of $\BBB$ satisfies
\begin{equation}\label{eq:bb-occupation}
  \mathbb{E}[\tau_{\BBB}(\Teps)]
  \leq K_b \cdot \exp((\Db - \alpha)/\eps) \cdot (1 + \delta_5),
\end{equation}
where $K_b > 0$ is a prefactor from the BEGK capacity formula for
$\BBB$ exit times (analogous to $K_s$), and
$|\delta_5| \leq C_5 \exp(-c_5/\eps)$ with
$c_5 \geq \min(\Ds - \Db, \Delta_{\mathrm{sub}}')$ ($\Delta_{\mathrm{sub}}'$
is the sub-leading barrier for $\BBB$).

\medskip
\textbf{Derivation.}
The excursion structure works as follows:
\begin{enumerate}
\item The chain starts in $\Bstable$ and must first exit $\Bstable$
  before visiting $\BBB$. The probability that the chain reaches
  $\BBB$ before time $\Teps$ is bounded by
  \[
    \mathbb{P}(\tau_{\Bstable^c} < \Teps)
    = \mathbb{P}\!\left(\frac{\tau_{\Bstable^c}}
      {\mathbb{E}[\tau_{\Bstable^c}]}
      < \frac{\Teps}{\mathbb{E}[\tau_{\Bstable^c}]}\right).
  \]
  Now $\mathbb{E}[\tau_{\Bstable^c}] = K_s \exp(\Ds/\eps)(1+\delta_2)$
  and $\Teps = \exp((\Ds - \alpha)/\eps)$, so
  \[
    \frac{\Teps}{\mathbb{E}[\tau_{\Bstable^c}]}
    = \frac{\exp(-\alpha/\eps)}{K_s(1+\delta_2)}
    = O(\exp(-\alpha/\eps)) \to 0.
  \]
  By the exponential law \eqref{eq:exp-law},
  $\mathbb{P}(\tau_{\Bstable^c}/\mathbb{E}[\tau_{\Bstable^c}] < s)
  \to 1 - e^{-s} \sim s$ for small $s$. Therefore
  \begin{equation}\label{eq:excursion-prob}
    \mathbb{P}(\text{any BB excursion before } \Teps)
    \leq \frac{\Teps}{\mathbb{E}[\tau_{\Bstable^c}]}(1 + o(1))
    = O(\exp(-\alpha/\eps)).
  \end{equation}

\item Conditional on an excursion to $\BBB$ occurring, the time spent
  in $\BBB$ during one visit is
  $\mathbb{E}[\tau_{\BBB \to \BBB^c}] = K_b \exp(\Db/\eps)(1+\delta_5')$
  by applying BEGK Theorem~1.2 to $\BBB$ (same formula as
  Lemma~\ref{lem:residence} but for the shallower basin).

\item Therefore the expected total BB occupation is bounded by
  \[
    \mathbb{E}[\tau_{\BBB}(\Teps)]
    \leq \mathbb{P}(\text{excursion occurs})
    \cdot \mathbb{E}[\text{BB occupation} \mid \text{excursion}]
    \cdot (1 + o(1))
  \]
  \[
    \leq \frac{\Teps}{\mathbb{E}[\tau_{\Bstable^c}]}
    \cdot K_b \exp(\Db/\eps) \cdot (1 + O(\exp(-c_5/\eps)))
  \]
  \[
    = \frac{\exp((\Ds - \alpha)/\eps)}{K_s \exp(\Ds/\eps)}
    \cdot K_b \exp(\Db/\eps)
    \cdot (1 + O(\exp(-c_5/\eps)))
  \]
  \begin{equation}\label{eq:bb-expected}
    = \frac{K_b}{K_s} \exp((\Db - \alpha)/\eps)
    \cdot (1 + \delta_5).
  \end{equation}
\end{enumerate}

\medskip
\textbf{BB experiential measure upper bound.}
Since $\rho(p) \leq \rho_{\max} = H(B)/4$ for all distributions $p$,
\begin{equation}\label{eq:mu-bb}
  \mu_{\BBB} \leq \rho_{\max} \cdot \mathbb{E}[\tau_{\BBB}(\Teps)]
  \leq \rho_{\max} \cdot \frac{K_b}{K_s}
  \cdot \exp((\Db - \alpha)/\eps) \cdot (1 + \delta_5).
\end{equation}

\medskip
\textbf{Rate constant:}
\begin{equation}\label{eq:c5}
  c_5 \geq \min(\Ds - \Db, \Delta_{\mathrm{sub}}') > 0.
\end{equation}

\medskip
\textbf{Output type:}
Upper bound on $\mathbb{E}[\tau_{\BBB}(\Teps)]$ and hence on
$\mathbb{E}[\mu_{\BBB}]$ (expected BB experiential measure), with
multiplicative error $(1 + \delta_5)$ where
$\delta_5 = O(\exp(-c_5/\eps))$.

Note: this is an expected-value bound, not a high-probability bound.
The conversion to a high-probability statement is handled by
Lemma~\ref{lem:concentration}.
\end{lemma}
%----------------------------------------------------------------------

%----------------------------------------------------------------------
\begin{lemma}[Concentration of Occupation-Time Functionals]%
\label{lem:concentration}
\textbf{Source:} For the exit-time concentration route:
BEGK \cite{BEGK02} Theorem~1.4 (exponential law of exit times).
For the general large-deviation route:
Donsker--Varadhan \cite{DV75-83} (large deviations for occupation
measures of Markov processes); see also den~Hollander,
\emph{Large Deviations}, AMS (2000), Chapter~VIII.

\medskip
\textbf{Statement (exit-time concentration route --- sufficient for
Theorem~A).}

By the BEGK exponential law (Lemma~\ref{lem:residence},
eq.~\eqref{eq:exp-law}), the exit time $\tau_{\Bstable^c}$
concentrates around its mean. Since the experiential measure
$\mu_{\mathrm{stable}}$ is controlled by the occupation time in
$\Bstable$ and the QSD convergence (Lemma~\ref{lem:qsd}), the
following hold simultaneously with probability at least
$1 - 2e^{-\eta/2} - o(1)$ as $\eps \to 0$:

\begin{align}
  \mu_{\mathrm{stable}}
  &\geq c \cdot K_s \cdot (1-\eta) \exp(\Ds/\eps)
    \cdot (1 - \delta_4),
    \label{eq:conc-stable} \\
  \mu_{\BBB}
  &\leq \rho_{\max} \cdot \frac{K_b}{K_s}
    \cdot \exp((\Db - \alpha)/\eps)
    \cdot (1 + \delta_5 + \delta_6),
    \label{eq:conc-bb}
\end{align}
where $\delta_6$ accounts for the deviation of $\tau_{\BBB}$ from its
expectation. For the BB occupation, since the excursion probability
itself is $O(\exp(-\alpha/\eps)) \to 0$, the random variable
$\tau_{\BBB}(\Teps)$ is zero with probability
$1 - O(\exp(-\alpha/\eps))$, and conditional on being nonzero,
the single-visit duration concentrates around its BEGK mean via
the exponential law applied to $\BBB$ exit.

Therefore $\delta_6 = O(\exp(-c_6/\eps))$ with
$c_6 = \alpha > 0$.

\medskip
\textbf{Statement (Donsker--Varadhan route --- for generality).}

For the time-averaged empirical measure
$L_T = \frac{1}{T}\int_0^T \mathbf{1}_{X_t \in A}\, dt$
of an ergodic continuous-time Markov chain with generator $Q$ on
a finite state space, the Donsker--Varadhan large deviation principle
gives:
\[
  \mathbb{P}(|L_T - \pi(A)| > \eta)
  \leq \exp(-T \cdot I_{\mathrm{DV}}(\eta))
\]
where $I_{\mathrm{DV}}(\eta) > 0$ for $\eta > 0$ and $\pi(A)$ is
the stationary occupation fraction. Since
$T = \Teps = \exp((\Ds - \alpha)/\eps)$ is exponentially large,
this gives super-exponentially tight concentration.

For weighted functionals $\frac{1}{T}\int_0^T f(X_t)\, dt$ with
bounded $f$ (as is the case for $\rho$ on finite state spaces, since
$0 \leq \rho \leq H(B)/4$), the contraction principle applied to
the DV rate function yields concentration with rate $I_f(\eta) > 0$.

\medskip
\textbf{Error term (exit-time route):}
\begin{equation}\label{eq:delta6}
  \delta_6 = O(\exp(-\alpha/\eps)),
  \qquad c_6 = \alpha.
\end{equation}

\medskip
\textbf{Note on the DV extension.}
The standard DV theorem addresses the unweighted empirical measure.
To handle the $\rho$-weighted functional
$\mu/T = \frac{1}{T}\int_0^T \rho(p_t)\, dt$, one applies the
contraction principle: since $\rho$ is a bounded continuous functional
on the space of distributions over a finite state space, the map
$L_T \mapsto \rho(L_T)$ is Lipschitz (with Lipschitz constant
$L_\rho < \infty$), and the rate function for the pushed-forward
functional is
$I_\rho(\eta) \geq I_{\mathrm{DV}}(\eta / L_\rho) > 0$.

However, for the present application (Theorem~A), the simpler
exit-time concentration route via BEGK Theorem~1.4 suffices and
avoids this extension entirely.

\medskip
\textbf{Output type:}
High-probability bounds on $\mu_{\mathrm{stable}}$ (lower) and
$\mu_{\BBB}$ (upper), converting the expected-value results of
Lemmas~\ref{lem:stable-lower} and~\ref{lem:bb-upper} into
probabilistic statements.
\end{lemma}
%----------------------------------------------------------------------

%----------------------------------------------------------------------
\begin{lemma}[Ratio Assembly]\label{lem:ratio}
\textbf{Source:} Direct combination of
Lemma~\ref{lem:stable-lower} (lower bound on $\mu_{\mathrm{stable}}$),
Lemma~\ref{lem:bb-upper} (upper bound on $\mu_{\BBB}$),
and Lemma~\ref{lem:concentration} (concentration).

\medskip
\textbf{Statement.}
Under the standing assumptions, as $\eps \to 0$:
\begin{equation}\label{eq:ratio}
  \frac{\mu_{\BBB}}{\mu_{\mathrm{stable}}}
  \leq C \cdot \exp\!\left(
    -\frac{\Ds - \Db - \alpha}{\eps}
  \right)
  \cdot (1 + \delta_7),
\end{equation}
where the prefactor is
\begin{equation}\label{eq:C-prefactor}
  C = \frac{\rho_{\max}}{c} \cdot \frac{K_b}{K_s^2},
\end{equation}
and
\begin{equation}\label{eq:delta7}
  |\delta_7| \leq C_7 \exp(-c_7/\eps),
  \qquad
  c_7 = \min(c_2, c_4, c_5, c_6)
  = \min(\Ds - \Db, \Ds, \alpha, \Delta_{\mathrm{sub}}, \Delta_{\mathrm{sub}}').
\end{equation}

\medskip
\textbf{Derivation.}
Using the high-probability bounds from Lemma~\ref{lem:concentration}:
\begin{align*}
  \frac{\mu_{\BBB}}{\mu_{\mathrm{stable}}}
  &\leq \frac{\rho_{\max} \cdot \frac{K_b}{K_s}
    \exp((\Db - \alpha)/\eps)(1 + \delta_5 + \delta_6)}
  {c \cdot K_s \cdot (1 - \eta) \exp(\Ds/\eps)(1 - \delta_4)} \\
  &= \frac{\rho_{\max} K_b}{c K_s^2 (1-\eta)}
  \cdot \exp\!\left(\frac{\Db - \alpha - \Ds}{\eps}\right)
  \cdot \frac{(1 + \delta_5 + \delta_6)}{(1 - \delta_4)}.
\end{align*}

Taking $\eta \to 0$ (which is valid since the probability
$1 - 2e^{-\eta/2}$ can be made arbitrarily close to $1 - 2/e
\approx 0.264$ and then improved by taking $\eta$ small),
and collecting the error terms:
\[
  \frac{(1 + \delta_5 + \delta_6)}{(1 - \delta_4)}
  = 1 + \delta_5 + \delta_6 + \delta_4 + O(\delta_i \delta_j)
  = 1 + O(\exp(-c_7/\eps)),
\]
since each $\delta_i = O(\exp(-c_i/\eps))$ with $c_i > 0$, and
$c_7 = \min_i c_i > 0$.

\medskip
\textbf{Verification that $C$ is $O(1)$ in $\eps$:}
\begin{itemize}
\item $\rho_{\max} = H(B)/4 = \ln|B|/4$: depends on $|B|$, not $\eps$.
\item $c = \rho(\nu_\QSD) > 0$: the QSD on $\Bstable$ has nonzero
  mutual information by standing assumption; $c$ depends on the
  basin structure but not on $\eps$ (the QSD converges to the
  ground-state conditional distribution as $\eps \to 0$).
\item $K_s, K_b$: BEGK prefactors are rational functions of the
  rate prefactors $r(x,y)$, the energy values $E(x)$, and the
  partition function $Z_\eps$. For Metropolis dynamics on finite
  state spaces, $K_s$ and $K_b$ are $O(1)$ in $\eps$ (they involve
  saddle-point curvatures and a finite number of rates, all bounded).
  See \cite{BdH15} Theorem~7.8 for explicit formulas.
\end{itemize}
Therefore $C = O(1)$ in $\eps$, confirming the exponential bound
is not spoiled by prefactor growth.

\medskip
\textbf{Critical check: error rate vs.\ exponent.}
The decay rate in the main term is $(\Ds - \Db - \alpha)/\eps$.
The error rate is $c_7/\eps$. We need $c_7 > 0$ for the error to
be negligible. We have:
\begin{align*}
  c_2 &\geq \Ds - \Db > 0, \\
  c_4 &= \Ds > 0, \\
  c_5 &\geq \Ds - \Db > 0, \\
  c_6 &= \alpha > 0.
\end{align*}
All are strictly positive. Moreover,
$c_7 \geq \min(\Ds - \Db, \alpha) > 0$
(since $\alpha > 0$ by assumption and $\Ds > \Db$).

Thus the error product remains $1 + O(\exp(-c_7/\eps))$ with
$c_7 > 0$, and the exponential decay $\exp(-(\Ds - \Db - \alpha)/\eps)
\to 0$ dominates.

\medskip
\textbf{Output type:}
The Theorem~A bound:
$\mu_{\BBB} / \mu_{\mathrm{stable}} \leq C \exp(-(\Ds - \Db - \alpha)/\eps)
(1 + O(\exp(-c_7/\eps)))$,
matching the statement in Section~7.2 of the draft.
\end{lemma}
%----------------------------------------------------------------------

%======================================================================
\section{Typed Dependency Graph}
\label{sec:dependency-graph}
%======================================================================

The seven lemmas form a directed acyclic graph (DAG).
Each edge is annotated with the mathematical object that flows from
source to target and the type of that object.

\begin{verbatim}
L1 (Basin Partition -- exact)
 |-- L2 (Residence Time -- expectation + concentration)
 |    |-- L4 (Stable Measure Lower Bound -- high-prob bound)
 |    +-- L5 (BB Occupation Upper Bound -- expected value)
 |-- L3 (QSD Convergence -- TV bound, rate gamma_D)
 |    +-- L4
 +-- L6 (Concentration -- converts E[.] to high-prob bounds)
      +-- L7 (Ratio Assembly -- Theorem A bound)
              ^
       L4, L5 -+
\end{verbatim}

\medskip
\noindent
\textbf{Topological order:}
$L1 \to (L2, L3, L6) \to (L4, L5) \to L7$.

\medskip
\noindent
\textbf{Edge table:}

\medskip
\begin{center}
\renewcommand{\arraystretch}{1.4}
\begin{tabular}{|c|c|p{4.2cm}|p{4.2cm}|p{3.0cm}|}
\hline
\textbf{Source} & \textbf{Target} & \textbf{Output (source)} & \textbf{Input (target)} & \textbf{Type match?} \\
\hline\hline

L1 & L2 &
  Basin $\Bstable$ with comm.\ height $\Ds$.
  \emph{Type: exact partition.} &
  A metastable set with known comm.\ height for BEGK capacity formula.
  \emph{Type: exact partition.} &
  \textbf{Yes.} Exact $\to$ exact. \\
\hline

L1 & L3 &
  Basin $\Bstable$ (as a subset of $\Omega$).
  \emph{Type: exact partition.} &
  A subset on which to define the killed chain.
  \emph{Type: exact partition.} &
  \textbf{Yes.} Exact $\to$ exact. \\
\hline

L1 & L6 &
  Basin partition and comm.\ heights.
  \emph{Type: exact partition.} &
  State space structure for DV / concentration argument.
  \emph{Type: exact partition.} &
  \textbf{Yes.} Exact $\to$ exact. \\
\hline

L2 & L4 &
  $\mathbb{E}[\tau_{\Bstable^c}] = K_s e^{\Ds/\eps}(1+\delta_2)$.
  \emph{Type: expectation with mult.\ error.}
  Also: exponential law giving concentration. &
  Residence time bound (needs at least a high-prob.\ bound on $\tau_{\Bstable^c}$, not just the mean).
  \emph{Type: high-probability bound.} &
  \textbf{Yes, via BEGK Thm~1.4.} The exponential law converts the expectation to a high-prob.\ bound. See \eqref{eq:exp-concentration}. \\
\hline

L3 & L4 &
  $\nu_\QSD$ with convergence rate $\gamma_D = O(1)$.
  \emph{Type: TV convergence bound.} &
  QSD to evaluate $\rho(\nu_\QSD) \geq c$.
  Convergence rate to bound transient error.
  \emph{Type: distribution + rate.} &
  \textbf{Yes.} TV convergence gives both the limiting distribution and the transient error bound needed for $\delta_4$. \\
\hline

L2 & L5 &
  $\mathbb{E}[\tau_{\Bstable^c}]$ (mean exit time, giving excursion probability $\Teps / \mathbb{E}[\tau]$).
  Also: BEGK for $\BBB$ exit time.
  \emph{Type: expectation.} &
  Excursion probability and conditional BB duration.
  \emph{Type: expectation $\to$ expected occupation.} &
  \textbf{Yes.} The excursion probability uses the ratio $\Teps / \mathbb{E}[\tau]$, which is exact up to $\delta_2$. The output is an expected-value bound, flagged for conversion by L6. \\
\hline

L4 & L7 &
  $\mu_{\mathrm{stable}} \geq c K_s e^{\Ds/\eps}(1-\delta_4)$.
  \emph{Type: lower bound on random variable.} &
  Denominator of the ratio.
  \emph{Type: lower bound (high-prob.).} &
  \textbf{Yes, after L6.} L6 converts L4's per-trajectory bound into a high-probability statement. \\
\hline

L5 & L7 &
  $\mathbb{E}[\mu_{\BBB}] \leq \rho_{\max} \frac{K_b}{K_s} e^{(\Db-\alpha)/\eps}(1+\delta_5)$.
  \emph{Type: expected-value upper bound.} &
  Numerator of the ratio.
  \emph{Type: high-probability upper bound.} &
  \textbf{Yes, after L6.} L6 converts the expected-value bound to a high-prob.\ bound via exit-time concentration. \\
\hline

L6 & L7 &
  High-probability bounds \eqref{eq:conc-stable}--\eqref{eq:conc-bb}.
  \emph{Type: high-probability bounds.} &
  Numerator and denominator as high-prob.\ bounds for division.
  \emph{Type: high-probability bounds.} &
  \textbf{Yes.} Both numerator (upper) and denominator (lower) are high-prob.\ bounds; their ratio gives a high-prob.\ upper bound on $\mu_{\BBB} / \mu_{\mathrm{stable}}$. \\
\hline
\end{tabular}
\end{center}

\medskip
\noindent
\textbf{DAG verification:}
\begin{itemize}
\item L1 has no incoming edges (root).
\item L7 has no outgoing edges (sink).
\item No cycles: all edges go from lower-numbered to higher-numbered
  lemmas (L1 $\to$ L2, L3, L6; L2 $\to$ L4, L5; L3 $\to$ L4;
  L4, L5, L6 $\to$ L7).
\item Every lemma's inputs are supplied by lemmas earlier in
  topological order.
\end{itemize}

\medskip
\noindent
\textbf{Type mismatch analysis:}

The critical potential mismatch identified in the plan (Pitfall~2:
expectations flowing where high-probability bounds are needed)
appears at two edges:

\begin{enumerate}
\item \textbf{L2 $\to$ L4:} L2 provides $\mathbb{E}[\tau_{\Bstable^c}]$
  (an expectation), but L4 needs a lower bound on the random variable
  $\tau_{\Bstable^c}$ (for the integral bound on $\mu_{\mathrm{stable}}$).

  \emph{Resolution:} BEGK Theorem~1.4 (\eqref{eq:exp-law}) provides
  $\tau / \mathbb{E}[\tau] \xrightarrow{d} \mathrm{Exp}(1)$, which gives
  high-probability lower bounds on $\tau$. This is a genuine
  conversion step, not an implicit assumption.

\item \textbf{L5 $\to$ L7:} L5 provides an expected-value upper bound
  on $\mu_{\BBB}$, but L7 divides two random quantities and needs
  high-probability control.

  \emph{Resolution:} Lemma~\ref{lem:concentration} (L6) explicitly
  handles this conversion. The BB occupation is zero with probability
  $1 - O(\exp(-\alpha/\eps))$, so the high-probability bound is
  actually simpler than the expected-value bound.
\end{enumerate}

Both type mismatches are resolved by cited results. No hidden
type conversions remain.

\medskip
\noindent
\textbf{L7 output vs.\ Theorem~A statement:}

Lemma~\ref{lem:ratio} produces
\[
  \frac{\mu_{\BBB}}{\mu_{\mathrm{stable}}}
  \leq C \exp\!\left(-\frac{\Ds - \Db - \alpha}{\eps}\right)
  (1 + O(\exp(-c_7/\eps))),
\]
which matches the Theorem~A statement in Section~7.2 of the draft:
\[
  \mu_{\BBB} / \mu_{\mathrm{stable}}
  \leq C \cdot \exp(-(\Ds - \Db - \alpha)/\eps).
\]
The explicit error term $(1 + O(\exp(-c_7/\eps)))$ with
$c_7 = \min(\Ds - \Db, \alpha) > 0$ is the refinement
contributed by this plan.

%======================================================================
\section{Error Rate Summary}
\label{sec:error-summary}
%======================================================================

\begin{center}
\renewcommand{\arraystretch}{1.3}
\begin{tabular}{|c|l|c|l|}
\hline
\textbf{Lemma} & \textbf{Error term} & \textbf{Rate $c_i$}
  & \textbf{Source} \\
\hline\hline
L1 & (none -- exact) & --- & FW cycle hierarchy \\
\hline
L2 & $\delta_2 = O(e^{-c_2/\eps})$
  & $c_2 \geq \Ds - \Db$
  & BEGK Thm 1.2 \\
\hline
L3 & $\delta_3(t) = C_3 e^{-\gamma_D t}$
  & $\gamma_D = O(1)$
  & CV16 Thm 2.1 \\
\hline
L4 & $\delta_4 = O(e^{-\Ds/\eps})$
  & $c_4 = \Ds$
  & Combination of L2, L3 \\
\hline
L5 & $\delta_5 = O(e^{-c_5/\eps})$
  & $c_5 \geq \Ds - \Db$
  & BEGK Thm 1.2 for $\BBB$ \\
\hline
L6 & $\delta_6 = O(e^{-\alpha/\eps})$
  & $c_6 = \alpha$
  & BEGK Thm 1.4 \\
\hline
L7 & $\delta_7 = O(e^{-c_7/\eps})$
  & $c_7 = \min(\Ds - \Db, \alpha)$
  & Composition \\
\hline
\end{tabular}
\end{center}

\medskip
\noindent
\textbf{Critical check:} All $c_i > 0$ (since $\Ds > \Db$ and
$\alpha > 0$). The final composite error rate
$c_7 = \min(\Ds - \Db, \alpha) > 0$ ensures the error product
does not spoil the exponential decay $\exp(-(\Ds - \Db - \alpha)/\eps)$.

Since $\alpha < \Ds - \Db$ (standing assumption), we have
$c_7 = \alpha$ in the generic case, and the error term decays
at the same rate as the observation-horizon parameter. This is
tight: choosing $\alpha$ closer to $\Ds - \Db$ sharpens the
main bound but worsens the error term, while choosing $\alpha$
small worsens the main bound but the error is negligible. The
optimal trade-off is a separate question for Plan~02.

%======================================================================
% Bibliography (inline for self-containment)
%======================================================================

\begin{thebibliography}{99}

\bibitem{BEGK02}
A.~Bovier, M.~Eckhoff, V.~Gayrard, and M.~Klein,
\emph{Metastability in reversible diffusion processes I:
Sharp asymptotics for capacities and exit times},
Comm.\ Math.\ Phys.\ \textbf{228} (2002) 219--255.

\bibitem{BdH15}
A.~Bovier and F.~den~Hollander,
\emph{Metastability: A Potential-Theoretic Approach},
Grundlehren der mathematischen Wissenschaften \textbf{351},
Springer (2015).

\bibitem{CV16}
N.~Champagnat and D.~Villemonais,
\emph{Exponential convergence to quasi-stationary distribution
and $Q$-process},
Probab.\ Theory Relat.\ Fields \textbf{164} (2016) 243--283.

\bibitem{DV75-83}
M.~D.~Donsker and S.~R.~S.~Varadhan,
\emph{Asymptotic evaluation of certain Markov process expectations
for large time},
Comm.\ Pure Appl.\ Math.\ I: \textbf{28} (1975) 1--47;
II: \textbf{28} (1975) 279--301;
III: \textbf{29} (1976) 389--461;
IV: \textbf{36} (1983) 183--212.

\bibitem{FW12}
M.~I.~Freidlin and A.~D.~Wentzell,
\emph{Random Perturbations of Dynamical Systems},
3rd ed., Grundlehren der mathematischen Wissenschaften \textbf{260},
Springer (2012).

\end{thebibliography}

\end{document}
